% SOURCE_AUTHORITY: sga5_fr_workpass.tex, lines 13399--13439 (LF-normalized unit).
% SOURCE_SHA256: 791F4EFFC5E02832D5D77ED03518C8156D6F07E4C8238B03545DB93D883FBB28
\section*{4. Generalidades sobre los términos locales}

Supóngase que $X'$ es liso sobre $k$ y, además, que el número de puntos fijos de $f$ es finito. Resulta, en particular, que el número de puntos fijos de $g^{-1}f^{\prime}$ es finito para cada $g \in G$.

Para cada punto fijo $x \in X$ de $f$, se introduce la función $\varphi$-central
\[
S_x(f',f,G,\varphi) : G \longrightarrow \bQ \quad ,
\]
con valores racionales, definida por la fórmula:

\begin{equation}\tag{4.1}
S_x(f',f,G,\varphi)(g) = \frac{1}{c^\varphi(g)} \sum_{x' \in {X'_x}^{g^{-1}f'}} (\nu_{x'}(g^{-1}f') - 1) \quad .
\end{equation}

Se han utilizado las notaciones siguientes:

\begin{equation}\tag{4.2}
c^\varphi(g) = \card G_g^\varphi \quad , \quad G_g^\varphi = \{ x \in G \mid \varphi(x)gx^{-1} = g \} \quad ,
\end{equation}

\begin{equation}\tag{4.3}
{X'_x}^{g^{-1}f'} = \{ x' \in X' \mid p(x') = x \; , \; (g^{-1}f^{\prime})(x') = x' \} \quad .
\end{equation}

\begin{equation}\tag{4.4}
\nu_{x'}(g^{-1}f') = \text{la multiplicidad del punto fijo } x' \text{ de } g^{-1}f^{\prime} \quad .
\end{equation}

Por composición con las inclusiones canónicas
\[
\bQ \hookrightarrow \bQ_\ell
\]
la función $S_x(f',f,G,\varphi)$ engendra funciones $\varphi$-centrales $S_{x,\ell}(f',f,G,\varphi)$ sobre $G$ con valores en el cuerpo $\bQ_\ell$ de los números $\ell$-ádicos.

\medskip
\noindent
\textbf{Conjetura 4.5.} \emph{Para cada $\ell \neq \car(k)$, los valores de la función $S_{x,\ell}(f',f,G,\varphi)$ son enteros $\ell$-ádicos.}

\medskip

En lo que sigue se demostrará que la conjetura 4.5 es verdadera si la dimensión de $X$ es igual a 1.
